\documentclass{article}
\usepackage{amsmath}
\usepackage{amssymb}
\usepackage{graphicx}
\usepackage{subcaption}
\usepackage{braket}
\usepackage[left=1cm,right=1cm,top=1cm,bottom=2cm]{geometry}
\title{Homework 4}
\author{Kirstin Horvat - 301579693}
\begin{document}
\maketitle

\section*{Question 1a}
Let \(A = \begin{bmatrix} a & b \\ c & d \end{bmatrix}\) and \(B = \begin{bmatrix} e & f \\ g & h \end{bmatrix}\) \\
\begin{paragraph}
The hermitian operator is the complex conjugate transpose of a matrix. An example below shows what happens when the complex conjugate and transpose is applied to a matrix.
\end{paragraph}

\begin{align*}
\begin{split}
    A^* &= \begin{bmatrix} a^* & b^* \\ c^* & d^* \end{bmatrix} A^T = \begin{bmatrix} a & c \\ b & d \end{bmatrix} \\\\
\end{split}
\end{align*}

\begin{align*}
    \begin{split}
        A^{\dagger} &= \begin{bmatrix} a^* & c^* \\ b^* & d^* \end{bmatrix} \\ 
        B^{\dagger} &= \begin{bmatrix} e^* & g^* \\ f^* & h^* \end{bmatrix} \\\\
    \end{split}
\end{align*}

\begin{align*}
\begin{split}
    B^{\dagger}A^{\dagger} &= \begin{bmatrix} e^* & g^* \\ f^* & h^* \end{bmatrix}\begin{bmatrix} a^* & c^* \\ b^* & d^* \end{bmatrix} \\
    &= \begin{bmatrix} e^* a^* + g^* b^* & e^* c^* + g^* d^* \\ f^* a^* + h^* b^* & f^* c^* + h^* d^* \end{bmatrix} \\
    &= \begin{bmatrix} (ea + gb)^* & (ec + gd)^* \\ (fa + hb)^* & (fc + hd)^* \end{bmatrix} \\\\
\end{split}
\end{align*}


\begin{align*}
\begin{split}
   (AB) &= \begin{bmatrix} a & b \\ c & d \end{bmatrix} \begin{bmatrix} e & f \\ g & h \end{bmatrix} \\
   &= \begin{bmatrix} ae + bg & af + bh \\ ce + dg & cf + dh \end{bmatrix} \\
   (AB)^{\dagger} &= \begin{bmatrix} (ae + bg)^* & (ce + dg)^* \\ (af + bh)^* & (cf + dh)^* \end{bmatrix} \\
\end{split}
\end{align*}

\begin{align*}
    \begin{split}
        (AB)^{\dagger} = B^{\dagger}A^{\dagger} = \begin{bmatrix} (ae + bg)^* & (ce + dg)^* \\ (af + bh)^* & (cf + dh)^* \end{bmatrix} = \begin{bmatrix} (ea + gb)^* & (ec + gd)^* \\ (fa + hb)^* & (fc + hd)^* \end{bmatrix} \\\\
    \end{split}
\end{align*}

\section*{Question 1b}
\begin{align*}
    \begin{split}
        [A, B] &= AB - BA \\
        &= 
        \begin{bmatrix}
            a & b \\ c & d
        \end{bmatrix}
        \begin{bmatrix}
            e & f \\ g & h
        \end{bmatrix} -
        \begin{bmatrix}
            e & f \\ g & h
        \end{bmatrix}
        \begin{bmatrix}
            a & b \\ c & d
        \end{bmatrix} \\
        &= 
        \begin{bmatrix} 
            ae + bg & af + bh \\ 
            ce + dg & cf + dh 
        \end{bmatrix} - 
        \begin{bmatrix} 
            ea + fc & eb + fd \\ 
            ga + gc & gb + hd 
        \end{bmatrix} \\
        &= \begin{bmatrix}
            ae + bg - ea - fc & af + bh - eb - fd \\
            ce + dg - ga - gc & cf + dh - gb - hd
        \end{bmatrix} \\
        &= \begin{bmatrix}
            bg - fc & af + bh - eb - fd \\
            ce + dg - ga - gc & cf - gb
        \end{bmatrix} \\
        [A, B]^{\dagger}
        &=    
        \begin{bmatrix} 
            (bg - fc)^* &  (ce + dg - ga - gc)^* \\
            (af + bh - eb - fd)^* & (cf - gb)^*
        \end{bmatrix} \\\\
    \end{split}
\end{align*}

\begin{align*}
\begin{split}
    [B^{\dagger}, A^{\dagger}] &= B^{\dagger}A^{\dagger} - A^{\dagger}B^{\dagger} \\
    &= 
    \begin{bmatrix}
        e^* & g^* \\ f^* & h^* \\
    \end{bmatrix}
    \begin{bmatrix}
        a^* & c^* \\ b^* & d^* \\
    \end{bmatrix}
    - \begin{bmatrix}
        a^* & c^* \\ b^* & d^* \\   
    \end{bmatrix}
    \begin{bmatrix}
        e^* & g^* \\ f^* & h^* \\
    \end{bmatrix} \\
    &=
    \begin{bmatrix}
        e^* a^* + g^* b^* & e^* c^* + g^* d^* \\
        f^* a^* + h^* b^* & f^* c^* + h^* d^*
    \end{bmatrix}
    - \begin{bmatrix}
        a^* e^* + c^* f^* & a^* g^* + c^* h^* \\
        b^* e^* + d^* f^* & b^* g^* + d^* h^*
    \end{bmatrix} \\
    &=
    \begin{bmatrix}
       g^* b^* - c^* f^* & e^* c^* + g^* d^* -a^* g^* - c^* h^* \\
        f^* a^* + h^* b^* - b^* e^* - d^* f^* & f^* c^* - b^* g^* 
    \end{bmatrix} \\
    &= 
    \begin{bmatrix}
        (gb - cf)^* & (ec + gd - ag - ch)^* \\
        (fa + hb - be - df)^* & (fc - bg)^*
    \end{bmatrix} \\
\end{split}
\end{align*}

\begin{align*}
    \begin{split}
        [A, B]^{\dagger} &= [B^{\dagger}, A^{\dagger}] \\
        &= \begin{bmatrix} 
            (bg - fc)^* &  (ce + dg - ga - gc)^* \\
            (af + bh - eb - fd)^* & (cf - gb)^*
        \end{bmatrix} = \begin{bmatrix}
            (gb - cf)^* & (ec + gd - ag - ch)^* \\
            (fa + hb - be - df)^* & (fc - bg)^*
        \end{bmatrix} \\
    \end{split}
\end{align*}

\section*{Question 1c}
\text{Pauli matrices:}
\begin{align*}
    \begin{split}
        \sigma_x =
        \begin{bmatrix}
            0 & 1 \\
            1 & 0 \\
        \end{bmatrix}
    \end{split}
\end{align*}
\begin{align*}
    \begin{split}
        \sigma_y =
        \begin{bmatrix}
            0 & -i \\
            i & 0 \\
        \end{bmatrix}
    \end{split}
\end{align*}
\begin{align*}
    \begin{split}
        \sigma_z =
        \begin{bmatrix}
            1 & 0 \\
            0 & -1 \\
        \end{bmatrix}
    \end{split}
\end{align*}

\begin{paragraph}
    Since the pauli matrices are antisymmetric where $[\sigma_x, \sigma_y] = -[\sigma_y, \sigma_x]$, there are only 3 commutators
\end{paragraph}

\begin{align*}
    \begin{split}
    [\sigma_x, \sigma_y] &=
    \begin{bmatrix}
        0 & 1 \\
        1 & 0
    \end{bmatrix}
    \begin{bmatrix}
        0 & -i \\
        i & 0
    \end{bmatrix}
    - \begin{bmatrix}
        0 & -i \\
        i & 0
    \end{bmatrix}
    \begin{bmatrix}
        0 & 1 \\
        1 & 0
    \end{bmatrix} \\
    &= 
    \begin{bmatrix}
        i & 0 \\
        0 & -i
    \end{bmatrix}
    - \begin{bmatrix}
        -i & 0 \\
        0 & i
    \end{bmatrix} \\
    &= \begin{bmatrix}
        i + i & 0 \\
        0 & -i - i
    \end{bmatrix} \\
    &= 2i \begin{bmatrix}
        1 & 0 \\
        0 & -1
    \end{bmatrix} \\
    &= 2i\sigma_z \\
    \end{split}
\end{align*}

\begin{align*}
    \begin{split}
    [\sigma_z, \sigma_y] &=
    \begin{bmatrix}
        1 & 0 \\
        0 & -1
    \end{bmatrix}
    \begin{bmatrix}
        0 & -i \\
        i & 0
    \end{bmatrix}
    - \begin{bmatrix}
        0 & -i \\
        i & 0
    \end{bmatrix}
    \begin{bmatrix}
        1 & 0 \\
        0 & -1
    \end{bmatrix} \\
    &= 
    \begin{bmatrix}
        0 & -i \\
        -i & 0
    \end{bmatrix}
    - \begin{bmatrix}
        0 & i \\
        i & 0
    \end{bmatrix} \\
    &= \begin{bmatrix}
        0 & -i - i \\
        -i - i & 0
    \end{bmatrix} \\
    &= -2i \begin{bmatrix}
        0 & 1 \\
        1 & 0
    \end{bmatrix} \\
    &= -2i\sigma_x \\
    [\sigma_y, \sigma_z] &= -[\sigma_z, \sigma_y] = 2i\sigma_x \\
    \end{split}
\end{align*}


\text{Note: $i * i = -1$}
\begin{align*}
    \begin{split}
    [\sigma_z, \sigma_x] &=
    \begin{bmatrix}
        1 & 0 \\
        0 & -1
    \end{bmatrix}
    \begin{bmatrix}
        0 & 1 \\
        1 & 0
    \end{bmatrix}
    - \begin{bmatrix}
        0 & 1 \\
        1 & 0
    \end{bmatrix}
    \begin{bmatrix}
        1 & 0 \\
        0 & -1
    \end{bmatrix} \\
    &= 
    \begin{bmatrix}
        0 & 1 \\
        -1 & 0
    \end{bmatrix}
    - \begin{bmatrix}
        0 & -1\\
        1 & 0
    \end{bmatrix} \\
    &= \begin{bmatrix}
        0 & 1 + 1 \\
        -1 - 1 & 0
    \end{bmatrix} \\
    &= 2i \begin{bmatrix}
        0 & 1 \\
        -1 & 0
    \end{bmatrix} \\
    &= 2i \begin{bmatrix}
        0 & -i \\   
        i & 0
    \end{bmatrix} \\
    &= 2i\sigma_y \\
    \end{split}
\end{align*}

\newpage

\section*{Question 2a}
\begin{paragraph}
    Using an example of a projection into the positive x direction
\end{paragraph}
\begin{align*}
    \begin{split}
        \hat{P}_+ &= \ket{+}\bra{+} \\
        &= \begin{bmatrix} 1 \\ 0 \end{bmatrix}  \begin{bmatrix} 1 & 0 \end{bmatrix} \\
        &= \begin{bmatrix} 
            1 & 0 \\ 
            0 & 0
        \end{bmatrix} \\
        det(\hat{P}_+ - \lambda I) &= \begin{bmatrix}
            1 - \lambda & 0 \\
            0 & -\lambda
        \end{bmatrix} \\
        &= -\lambda (1 - \lambda) \\
        \lambda_1 = 0, \lambda_2 = 1
    \end{split}
\end{align*}
\section*{Question 2b}
Eigenvalues for the projection operator for the spin up state
\begin{align*}
    \begin{split}
        P_\uparrow &= \begin{bmatrix} 1 & 0 \\ 0 & 0 \end{bmatrix} \\
        det(P_\uparrow - \lambda I) &= \begin{bmatrix} 1 - \lambda & 0 \\ 0 & -\lambda \end{bmatrix} \\
        &= -\lambda (1 - \lambda) \\
        \lambda_1 &= 0 \\
        \lambda_2 &= 1 \\
    \end{split}
\end{align*}
Eigenstates (eigenvectors) for the projection operator for the spin up state
\begin{align*}
    \begin{split}
        \lambda = 1: \\
        &= \begin{bmatrix}
            1 - 1  & 0 \\
            0 & -1
        \end{bmatrix} \\
        &= \begin{bmatrix}
            0 & 0 \\
            0 & -1
        \end{bmatrix} \\\\
        \lambda = 0: \\
        &= \begin{bmatrix}
            1 & 0 \\
            0 & 0
        \end{bmatrix} \\\\
    \end{split}
\end{align*}
Eigenvalues of the projection operator for the spin down state
\begin{align*}
    \begin{split}
        P_\downarrow &= \begin{bmatrix} 0 & 0 \\ 0 & 1 \end{bmatrix} \\
        det(P_\downarrow - \lambda I) &= \begin{bmatrix} -\lambda & 0 \\ 0 & 1 - \lambda \end{bmatrix} \\
        &= -\lambda (1 - \lambda) \\
        \lambda_1 &= 0 \\
        \lambda_2 &= 1 \\
    \end{split}
\end{align*}
Eigenstates (eigenvectors) for the projection operator for the spin down state
\begin{align*}
    \begin{split}
        \lambda = 1: \\
        &= \begin{bmatrix}
            -1 & 0 \\
            0 & 1 - 1
        \end{bmatrix} \\
        &= \begin{bmatrix}
            -1 & 0 \\
            0 & 0
        \end{bmatrix} \\\\
        \lambda = 0: \\
        &= \begin{bmatrix}
            0 & 0 \\
            0 & 1
        \end{bmatrix} \\\\
    \end{split}
\end{align*}

\section*{Question 2c}

\begin{align*}
    \begin{split}
    \end{split}
\end{align*}


\end{document}

